\documentclass{article}
\usepackage{amsmath,amssymb,amsthm}
\newtheorem{theorem}{Theorem}
\begin{document}
\title{A064: exclusive-controller lift of priced prefix reserves}
\author{A064, B13}\date{14 October 2026}\maketitle
Entropy is in bits. Let a finite uniform physical controller select independent
Boolean cores $u_j$ on disjoint data cubes, or constant signs. Merge selectors
for the same physical core. Let $C_j$ be the exclusive indicators, $q_j=EC_j$,
$\ell_j=\log_2(1/q_j)$, with zero-$q_j$ terms zero. Controls are unqueried
during a finite adaptive data-only prefix. Physical histories retain fair
independent unrevealed coordinates. Every core follows a prescribed finite
continuation tree with a nonnegative reserve $r_j$ satisfying
\[
D_j\le C d_j+r_j-E r_j(\text{child}),\qquad C\ge0.
\]
Each $r_j$ depends only on its own core history and any initially fixed independent seed.
These core reserves are assumptions or independently supplied certificates,
not universal FEI conclusions. Interleaving respects those trees but can be
history-dependent and stopped at any reached prefix.

\begin{theorem}[Initial resource, drift, terminal boundary]
Put $\alpha=C+11/2$ and explicitly define
\[
R_t=\sum_j q_jr_j(u_{j,t})+\sum_j q_j\ell_j\operatorname{Var}(u_{j,t}).
\]
Then $R_t\ge0$ and at every actual data query
\[
D_F\le\alpha d_F+R_t-E[R_{t+1}\mid\mathcal F_t].
\]
Its initial resource satisfies
\[
R_0\le\sum_jq_jr_j(u_j)+2I_{\rm ctrl}(F).
\]
At a stopped prefix $\sigma$, its actual value is the displayed sum over the
remaining cores. In particular
\[
E R_\sigma\le R_0+C\Pi_{I,\rm data},
\]
where $\Pi_{I,\rm data}$ is the same prefix's consumed physical data influence.
If every core is exhausted to a constant and the supplied core reserves have
zero constant leaves, $R_\sigma=0$ exactly. The Fourier entropy of the remaining
control function $F_\sigma$ is still retained.
\end{theorem}
\begin{proof}
At a data query of core $j$, write $m=Eu_j$, $\beta=\widehat u_j(i)$,
$t=\beta^2$, $v=1-m^2$, and $d=\operatorname{Inf}_i(u_j)$. The exact
signed sector identity gives
\[
D_F=q_j[D_j-p(m,\beta)]+M(t)+J,
\quad d_F=q_jd.
\]
Here $\phi(s)=s\log_2(1/s)$ and
$p(a,b)=\phi(a^2)+\phi(b^2)-E_\pm\phi((a\pm b)^2)$.
It is nonnegative: with $\Phi(A,B)=\phi(A)+\phi(B)-\phi(A+B)$,
\[
p(a,b)=\Phi(a^2,b^2)+\phi(a^2+b^2)-E_\pm\phi((a\pm b)^2)\ge0.
\]
The first term is homogeneous binary entropy; the last difference is
nonnegative by concavity, including zero endpoints.

The independently reconstructed selector envelope is
\[
M(t)\le q_jt[\ell_j+1/\ln2]+q_jt\log_2(1/t),\quad
0\le J\le2q_jt/\ln2.
\]
Indeed $M(v)\le h_2(q_jv)$, and the scaling envelope
$M(t)\le tM(v)/v+q_jt\log_2(v/t)$ cancels the variance logarithm.
Boolean derivatives give $|\beta|\le d$ and $t\le d$, while
$t\log_2(1/t)\le2d/(e\ln2)$. Thus
\[
M(t)+J\le\kappa q_jd+q_j\ell_jt,\qquad
\kappa=(3+2/e)/\ln2<11/2.
\]
The rational certificate uses $e>65/24$, $\ln2>56/81$ and
$(3+48/65)(81/56)<11/2$. No floating constant comparison is used.
The core reserve and the variance drift $v-Ev_{\rm child}=t$ prove
the claimed inequality. Other cores have unchanged conditional quantities.

For the initial resource, indicator edge isoperimetry gives
$I(C_j)\ge2q_j\ell_j$ in probability-of-change normalization.
At an original control edge between distinct states $j,k$, actual pivotal
probability is $(1-m_jm_k)/2$, while the charged variance sum is at most
$2(1-m_jm_k)$. Constants have variance zero. Summing each original edge once
gives $\sum_jv_jI(C_j)\le4I_{\rm ctrl}$ and hence
$\sum_jq_j\ell_jv_j\le2I_{\rm ctrl}$. Edge isoperimetry follows by
cube slice induction and $h_2(s)\ge2\min(s,1-s)$, so the influence units
are explicit. For the expected terminal bound, finite summation of the core
reserve inequalities gives $E\sum q_jr_j(\sigma)\le\sum q_jr_j(0)
+C\Pi_{I,\rm data}-\sum_jq_jE\sum D_j$. Variances decrease in expectation.
The discarded terms are nonnegative. All statements are for finite trees;
no unbounded optional stopping or dimension-limit exchange is invoked.
\end{proof}

\begin{theorem}[Snell domination and actual controller continuation]
Let $Q_\alpha$ be the least nonnegative reserve on a prescribed physical tree:
$Q_\alpha(v)=[D_v-\alpha d_v+E Q_\alpha(\text{child})]_+$, zero at its
specified leaves. On the above data-only tree, $Q_\alpha(F)\le R_0$.
For a complete-data frontier, choose a declared control continuation at each
actual $F_\sigma$ and let $b_\sigma$ dominate its $Q_\alpha$.
Then adding $A_t=E[b_\sigma\mid\mathcal F_t]$ yields a valid nonnegative
potential for the full grafted rule and
\[
Q_\alpha(F,\text{full rule})\le
\sum_jq_jr_j(u_j)+2I_{\rm ctrl}(F)+E b_\sigma.
\]
No linear bound on $E b_\sigma$ is assumed. If one is proved at rate $L$
against this same remaining control influence, the total controller account
has rate $2+L$, not $\max(2,L)$.
\end{theorem}
\begin{proof}
Backward induction gives minimality and domination by any nonnegative
supersolution with dominating leaves. $A$ is a nonnegative martingale,
so adding it does not change the data drift. At the complete-data boundary
$R=0$, hence $R+A=b_\sigma$ matches the actual continuation reserve.
This glues the declared rule without changing its frontier or reference.
\end{proof}

\paragraph{Arbitrary retained boundary.}
For any Boolean sign $h(z)$, choose two independent dictator cores and put
$F=((1+h)/2)y_0+((1-h)/2)y_1$. Data-first leaves are constant when
$y_0=y_1$ and $\pm h$ otherwise. Therefore
\[
E H(F_\sigma)=H(h)/2,\quad E I(F_\sigma)=I(h)/2=I_{\rm ctrl}(F),
\quad E Q_\alpha(F_\sigma,T)=Q_\alpha(h,T)/2
\]
for consistently sign-translated copies of one declared control tree $T$.
Each data core has zero reserve. A universal linear entropy payment of this
actual frontier would be FEI. The new prefix potential does not supply it.

\paragraph{Growing branch.}
For $q=1/4$ selecting $u=(1,1,1,-1)$ with constant fallback and $C=4$,
the core reserve is zero. The rarity potential is $3/8$ at the root and
$1/2$ on its nonconstant first-data child, a strict increase of $1/8$.
Its other child has value zero. Conditional expected drift, rather than
pathwise monotonicity, is what the theorem proves.
\end{document}
